\chapter{Implementation Details \& Infrastructure}
\label{ch:implementation_details}

Moving from architectural design to a functional system requires selecting a technology stack that matches the concurrency and networking demands of the Raft protocol. This chapter justifies the choice of Go and gRPC and explains the implementation of the cross-platform client ecosystem.

\section{Technology Stack: Why Go and gRPC}
The selection of the technology stack was driven by the need for high-performance concurrency management and efficient, typed inter-node communication.

\subsection{The Go Programming Language}
The Go programming language was selected for the backend implementation due to its first-class support for the Communicating Sequential Processes (CSP) model. This paradigm is particularly well-suited for the Raft algorithm, which relies heavily on asynchronous state transitions and event-driven logic \cite{GoConcurrency2017}.

Key features of Go utilized in the project include:
\begin{itemize}
    \item \textbf{Goroutines}: These lightweight threads allow every node to manage multiple background tasks—such as the election ticker, log application, and per-peer replication—without the overhead of traditional OS threads \cite{raft_go}.
    \item \textbf{Channels and Select}: Channels serve as the primary synchronization mechanism between the Raft core and the Replicated State Machine (RSM). The \texttt{select} statement allows a node to wait on multiple asynchronous events (e.g., election timeouts vs. heartbeat arrivals) in a non-blocking manner \cite{raft_go, rsm_go}.
    \item \textbf{Memory Safety and Tooling}: Go's strict typing and integrated race detector were vital during the debugging of complex concurrency bugs in the log replication logic.
\end{itemize}

\subsection{gRPC and Protocol Buffers}
For the communication layer, gRPC was chosen over REST or custom TCP protocols. gRPC provides a high-performance framework built on top of HTTP/2, offering several advantages for a distributed state machine:
\begin{itemize}
    \item \textbf{Binary Serialization}: Using Protocol Buffers (\textit{protobuf}) significantly reduces the payload size of replicated logs compared to JSON, which is critical for maintaining high throughput \cite{GoogleProtobuf}.
    \item \textbf{Strict Typing}: The service definitions in \texttt{raft.proto} and \texttt{kv.proto} ensure that all cluster members agree on the structure of the data, reducing the likelihood of serialization errors during term updates or snapshot transfers \cite{10, 11}.
    \item \textbf{Streaming Capabilities}: gRPC's support for streaming allows the system to potentially handle large snapshot transfers more efficiently than discrete request-response cycles.
\end{itemize}

\section{Cross-Platform Interoperability}
A major architectural goal was to ensure that the key-value store could be accessed by applications regardless of their native language. This is achieved by separating the Go-based consensus engine from the client-side library.

\subsection{Integrating the Java Clerk with the Go Backend}
The project includes a production-grade Java client implementation, the \texttt{Clerk.java}, which demonstrates the system's cross-platform interoperability. This integration relies on the gRPC framework to bridge the gap between the Go server and the Java Virtual Machine (JVM).

The integration architecture follows a shared-contract model:
\begin{enumerate}
    \item \textbf{Shared Protobuf Interface}: Both the Go server and the Java client use the same \texttt{.proto} files to generate their respective networking code \cite{10, 11}.
    \item \textbf{Consistent Logic}: The Java \texttt{Clerk} implements the exact same leader discovery and retry logic as the Go client. It interprets \texttt{Status\_ERR\_WRONG\_LEADER} responses and uses the \texttt{leader\_hint} to redirect requests \cite{24, 25}.
    \item \textbf{Asynchronous Programming Model}: While the Go backend uses goroutines, the Java client utilizes \texttt{CompletableFuture} and gRPC \texttt{FutureStub} to provide a non-blocking API to Java applications, maintaining high responsiveness \cite{24}.
\end{enumerate}

The success of this cross-language integration proves that the Raft implementation is not tied to a specific ecosystem.
By relying on industry-standard protocols, the system can support a diverse range of clients, from low-latency Go microservices to heavy-duty Java enterprise applications.

\section{Deployment \& Automation}
\label{sec:deployment_automation}
To ensure the system is reproducible across different environments and easy to manage at scale, we utilize containerization and automated configuration scripts. This approach allows for the rapid instantiation of clusters and provides a consistent environment for evaluation.

\subsection{Containerization with Docker}
Every node in the Raft cluster is packaged as a lightweight Docker container. This ensures that the runtime environment, including dependencies and networking configurations, remains identical across development, testing, and production stages.

The system utilizes a \textbf{multi-stage build} process defined in the \texttt{Dockerfile} to optimize the final image size and security:
\begin{itemize}
    \item \textbf{Build Stage}: A \texttt{golang:1.25-alpine} image is used as the base to compile the \texttt{kvserver} binary. This stage includes downloading dependencies and performing a static build to ensure the binary is self-contained.
    \item \textbf{Final Stage}: The resulting binary is copied into a minimal \texttt{alpine:latest} image. This stage only contains the executable and necessary CA certificates, significantly reducing the attack surface and deployment footprint.
\end{itemize}

The \texttt{ENTRYPOINT} of the container is set to execute the \texttt{kvserver} binary, allowing the orchestrator to pass node-specific configuration via command-line flags or environment variables.

\subsection{Cluster Orchestration with Docker Compose}
The orchestration of multiple nodes and the supporting monitoring stack is managed via Docker Compose. As seen in the generated \texttt{docker-compose.yml}, the cluster typically consists of five nodes (\texttt{node0} through \texttt{node4}) and a suite of observability services including Prometheus, Grafana, Loki, and Promtail.

Key architectural features of the Compose configuration include:
\begin{itemize}
    \item \textbf{Host Networking}: The services use \texttt{network\_mode: "host"} to minimize network latency and allow nodes to communicate via their specific ports as defined in the cluster configuration.
    \item \textbf{Volume Persistence}: Each node is assigned a dedicated volume (e.g., \texttt{node0\_data}) to store its Raft log and snapshots, ensuring state is preserved across container restarts.
    \item \textbf{Log Aggregation}: The \texttt{logs\_data} volume is shared between the Raft nodes and Promtail, allowing the monitoring stack to ingest and index structured logs in real-time.
\end{itemize}

\subsection{Configuration Automation via setup\_cluster.py}
To simplify the management of different cluster topologies and timeout settings, a Python-based automation script, \texttt{setup\_cluster.py}, was developed. This script acts as a configuration generator that translates a simple \texttt{cluster.json} file into a complete deployment package.

The automation script handles the following tasks:
\begin{enumerate}
    \item \textbf{Validation}: Ensures the cluster topology contains an odd number of nodes (at least 3) to satisfy Raft's majority requirement.
    \item \textbf{Compose Generation}: Dynamically creates the \texttt{docker-compose.yml} with the appropriate environment variables for Raft timeouts (e.g., \texttt{election\_min}, \texttt{heartbeat}) and gRPC backoff parameters.
    \item \textbf{Monitoring Sync}: Automatically generates the \texttt{prometheus.yml} target list based on the nodes' IDs and metrics ports, ensuring the dashboard is always synchronized with the active cluster.
\end{enumerate}

This automation allows for rapid experimentation with different network latencies and consensus parameters, which is essential for the performance evaluation conducted in the following chapter.